\section{Discussion}
\label{sec:evidence_scope}

\paragraph{Anticipation for embodied musical response}
Music-driven humanoid dance requires preparation:
weight transfer and coordinated articulation develop
before a visible accent.
Forecast music states provide prospective context for
these decisions using arrived audio alone.
Our framework connects a pretrained musical continuation
prior to robot-native motion through paired dance learning.
The relevant quantity is the useful musical lead at
commitment, which depends jointly on forecast content,
temporal alignment, and processing delay.
This connects musical anticipation to the timing
requirements of embodied movement.

\paragraph{Structured expression across commitments}
Kinematics-guided D+C supervision assigns discrete codes
responsibility for designated whole-body organization
while continuous residuals retain complementary variation.
These roles are defined through reconstruction objectives
and expressed through a shared decoder, allowing structure
and detail to remain coupled.
\cof{} extends this coupling across planning cycles:
the committed components determine both the latent
history and the reference boundary state used for
continuation.
This formulation connects the representation of movement
to the consequences of previous generation decisions,
making their consistency an explicit part of training.

\paragraph{Preserving expression through execution}
A robot's musical response depends on how generated
movement survives reference delivery and physical tracking.
The robot-native interface connects our planner to
fixed whole-body control, while paired evaluation
examines musical correspondence and motion quality
alongside tracking accuracy.
This makes the preservation of expressive movement
a system-level objective.
The current generator plans from committed references,
with measured-state feedback handled by SONIC.
Extending the planning context to include execution
deviations could enable adaptation to the robot's
actual posture and support state.
Forecast uncertainty also motivates adapting planning
and commitment horizons when musical transitions
become less predictable.

\section{Conclusion}

We presented \method{}, a framework for streaming
music-driven humanoid dance that integrates musical
anticipation, robot-native D+C generation, and
committed continuation.
A frozen music foundation model forecasts upcoming
musical context from arrived audio.
Kinematics-guided supervision organizes discrete
structure and complementary continuous detail,
while \cof{} trains the generator to continue from
jointly committed motion and its resulting reference
state.
The generated stream is delivered to a fixed
whole-body controller for humanoid execution.

% TODO: Insert the principal experimental finding on
% musical adaptation, component benefits, and preservation
% during robot execution.

By connecting musical foresight to structured motion
and successive commitments, \method{} provides
a unified approach to generating humanoid dance
in response to music as it unfolds.